\section{Ewaluacja}\label{sec:ewaluacja}

Porównanie sekwencji referencyjnej i wygenerowanej rozpoczyna się od dopasowania rekordów za pomocą algorytmu węgierskiego \parencite{kuhn1955hungarian}.
Macierz podobieństwa budowana jest na podstawie ważonej zgodności pól rekordu, przy czym na etapie samego dopasowania implementacja korzysta ze znormalizowanych łańcuchów znaków i miary podobieństwa \texttt{SequenceMatcher} z biblioteki \texttt{difflib} \parencite{python_difflib_2026}.
Takie rozwiązanie pozwala zestawiać ze sobą sekwencje różniące się długością i kolejnością wpisów, a jednocześnie zachować próg jakości wymagany do uznania dwóch rekordów za odpowiadające sobie.

Po wykonaniu dopasowania obliczane są dwa typy ocen.
W wariancie \textit{source} dla każdego pola raportowane są miary podobieństwa tekstowego: średnia, minimum i maksimum.
W tym etapie wykorzystywany jest scorer oparty na znormalizowanym podobieństwie edycyjnym typu Levenshteina \parencite{wagner1974string}.
W wariancie \textit{parsed} ten sam typ porównania służy do wyznaczenia poprawności klasyfikacyjnej po przyłożeniu ustalonego progu zgodności.
Nie jest to jednak klasyczna macierz pomyłek budowana dla niezależnych obserwacji i zamkniętego zbioru etykiet.
Kategorie \textit{true positive}, \textit{false positive}, \textit{false negative} i \textit{true negative} wyznaczane są tu po wcześniejszym dopasowaniu rekordów, a następnie przez sprawdzenie, czy dana wartość pola została odtworzona w stopniu wystarczająco zgodnym z referencją.

Za trafienie (\textit{true positive}) uznaje się sytuację, w której zarówno wartość referencyjna, jak i przewidywana istnieją, a ich podobieństwo przekracza ustalony próg.
Wpis wygenerowany przez model, któremu nie odpowiada żadna wartość referencyjna, lub którego wartość odbiega od referencji poniżej progu, liczony jest jako błąd nadmiarowy (\textit{false positive}).
Analogicznie, wpis referencyjny bez odpowiadającej mu predykcji lub z predykcją poniżej progu stanowi pominięcie (\textit{false negative}).
W szczególności przypadek, w którym model zwrócił wartość dla pola, ale nie spełnia ona progu zgodności względem referencji, traktowany jest jednocześnie jako błąd nadmiarowy i pominięcie, ponieważ system wygenerował odpowiedź, lecz nie odtworzył poprawnej wartości.
Na tej podstawie wyznaczane są precyzja, czułość, dokładność oraz ich średnia harmoniczna (miara F1), oddzielnie dla każdego pola i dla każdego schematyzmu.

Wyniki ewaluacji eksportowane są w postaci arkuszy kalkulacyjnych, zawierających dla każdego schematyzmu zestawienie metryk per pole wraz z wierszem podsumowującym.
Nazwy plików wynikowych zawierają identyfikator użytego modelu oraz znacznik czasowy, co pozwala na systematyczne porównywanie kolejnych przebiegów eksperymentalnych.
Interpretacja tych metryk wymaga jednak rozbicia na poszczególne pola, ponieważ dalsza analiza jakościowa pokazuje wyraźne różnice między relatywnie łatwiejszym numerem strony i materiałem budowlanym a znacznie trudniejszym zadaniem poprawnego odziedziczenia dekanatu.
